\section{Introduction}\label{sec:introduction}

Let a holomorphic bundle become strictly semistable as a Kähler polarization
approaches a wall. On the stable side its Hermitian--Einstein metric may
diverge in a fixed holomorphic frame. The limiting polystable bundle has
additional parallel endomorphisms, and these give rise to small eigenvalues
of the endomorphism Dolbeault Laplacian. When a stable factor is repeated,
the corresponding parallel endomorphisms form a matrix algebra. The small
spectrum therefore contains matrix directions that are absent from a
scalar graph model.

We determine the entire small spectrum for simple extensions of two
polystable layers, uniformly as positive spectral scales separate. The
finite operator is specified by the extension class and a quadratic
moment-map equation. The principal analytical issue is to compare this
independently specified balance with the actual Hermitian--Einstein metric.
An ordinary implicit-function estimate for the finite equation would divide
by its smallest positive eigenvalue. Such an estimate is not uniform near
a face of a stability chamber. We instead use curvature identities for
all multiplicity projections and a finite-dimensional degree argument.

\subsection{Geometric data and normalization}
Let $X$ be a compact connected complex manifold of dimension $d$, and let
$\omega_s$, for $s$ in a neighborhood of $0$ in $\R^k$, be a smooth family
of Kähler forms. Put $dV_s=\omega_s^d/d!$ and $V_s=\int_XdV_s$.
Fix pairwise nonisomorphic locally free bundles $F_v$, of ranks $r_v$,
which are stable and have the same slope with respect to $\omega_0$.
The index set is the disjoint union of nonempty sets $I$ and $J$.
Fix positive multiplicities $m_v$ and a simple holomorphic extension
\begin{equation}\label{eq:extension}
0\longrightarrow S=\bigoplus_{i\in I}F_i\otimes\C^{m_i}
\longrightarrow E\longrightarrow
U=\bigoplus_{j\in J}F_j\otimes\C^{m_j}\longrightarrow0.
\end{equation}
Simple means $H^0(X,\End E)=\C\id_E$. Repetitions are allowed within
either layer; the same stable isomorphism type does not occur in both.
At least one $m_v>1$ in the applications of interest.

After shrinking the parameter neighborhood, suppose the factors have
smoothly varying Hermitian--Einstein metrics $h_{v,s}$. This is the usual
local continuation of the stable factor metrics; it can also be taken as
an explicit hypothesis on the chosen family. Fix their scalar
normalizations. Write $i\Lambda_sF_{h_{v,s}}=\kappa_v(s)\id_{F_v}$ and
\begin{equation}\label{eq:kappa}
R_E=\sum_vr_vm_v,\qquad
\kappa_E(s)=\frac{1}{R_E}\sum_vr_vm_v\kappa_v(s),\qquad
b_v(s)=-r_v\bigl(\kappa_v(s)-\kappa_E(s)\bigr).
\end{equation}
Our curvature convention is $c_1(F)=[i\tr F_h/(2\pi)]$; thus
\begin{equation}\label{eq:degree}
\kappa(F)=\frac{2\pi}{\rk(F)V_s}
\int_Xc_1(F)\wedge\frac{\omega_s^{d-1}}{(d-1)!}.
\end{equation}
All slope comparisons have the same sign in this normalization.
The numbers $b_v$ satisfy $\sum_vm_vb_v=0$. They are coefficients in
the dual trace pairing, so their definition already includes the factor
ranks.

Choose fixed bases of $\Ext^1(F_j,F_i)$ and take their current harmonic
representatives $\beta_{ij,a,s}$. If $Z=(Z_{ij,a})$ is an array of
$m_i$ by $m_j$ matrices, set
\begin{equation}\label{eq:harmonic}
N_s(Z)_{ij}=\sum_a\beta_{ij,a,s}\otimes Z_{ij,a},\qquad
\dbar_{s,Z}=\dbar_G+N_s(Z),\qquad G=S\oplus U,
\end{equation}
and define $q_s(Z)=V_s^{-1}\norm{N_s(Z)}_{L^2}^2$ using
$h_{0,s}=\bigoplus_vh_{v,s}\otimes I_{m_v}$. We retain the actual
harmonic Gram matrices. Fixed cohomology coordinates need not be
orthonormal for varying $s$. The extension $E$ has fixed cohomology
matrices $Z^0$ in this notation.

The group $\mathcal G=\prod_v\mathrm{GL}(m_v,\C)$ acts by
$(g\cdot Z)_{ij,a}=g_iZ_{ij,a}g_j^{-1}$; its compact subgroup is
$K=\prod_v\mathrm U(m_v)$. Define the quadratic moment map by
\begin{equation}\label{eq:momentdefinition}
\sum_v\tr\bigl(Y_v\mu_s(Z)_v\bigr)
=\left.\frac{d}{dt}\right|_{t=0}q_s(e^{tY/2}\cdot Z)
\quad(Y_v=Y_v^*).
\end{equation}
In orthonormal arrow coordinates it is the sum of incoming $ZZ^*$ minus
outgoing $Z^*Z$. We call $Z$ \emph{balanced} at $s$ if
$\mu_s(Z)_v=b_v(s)I_{m_v}$. Finite balancing below concerns the orbit
$\mathcal G\cdot Z^0$, independently of the unknown HE metric.

A collection of subspaces $W_v\subset\C^{m_v}$ is a
\emph{subrepresentation} if $Z^0_{ij,a}W_j\subset W_i$ for every arrow.
It is proper and nonzero if it is neither the whole collection nor zero.
The representation is $b(s)$-stable when
\begin{equation}\label{eq:finite-stability}
\sum_vb_v(s)\dim W_v>0
\quad\text{for every proper nonzero subrepresentation }W.
\end{equation}
This is King's positive sign convention. The real-weight balancing theorem
then supplies a balanced $Z\in\mathcal G\cdot Z^0$, unique up to $K$.
The details and sign conversion are given in Section~\ref{sec:extensions}.

\subsection{The finite spectrum and the main theorem}
Consider the full complex vector space
\begin{equation}\label{eq:finite-space}
\HH=\left\{X=(X_v)\in\bigoplus_vM_{m_v}(\C):
\sum_vr_v\tr X_v=0\right\},\qquad
\norm X_r^2=\sum_vr_v\tr(X_v^*X_v).
\end{equation}
Its complex dimension is $N=\sum_vm_v^2-1$. For a balanced $Z$, define
\begin{equation}\label{eq:finite-form}
\mathcal Q_{s,Z}(X)=q_s([X,Z]),\qquad
[X,Z]_{ij,a}=X_iZ_{ij,a}-Z_{ij,a}X_j.
\end{equation}
Let $\nu_1(s)\le\cdots\le\nu_N(s)$ be its eigenvalues relative to
$\|\cdot\|_r^2$.
Simplicity makes these eigenvalues positive. Compact-unitary invariance
makes them independent of the choice of balanced representative.
They are defined from the geometry, extension orbit, and finite balance.

\begin{theorem}\label{thm:main}
Fix the geometric data above. There are a neighborhood $\mathcal U$ of
$0$ and constants $C,\gamma>0$ with the following properties.
For $s\in\mathcal U$, the bundle $E$ is slope-stable with respect to
$\omega_s$ if and only if its multiplicity representation is
$b(s)$-stable. At such a parameter let $Z$ be balanced and put
\begin{equation}\label{eq:sigma}
\sigma(s)=\sum_{i\in I}m_ib_i(s)=q_s(Z)>0.
\end{equation}
Let $h_{\HE,s}$ be the HE metric whose determinant is the product
determinant metric induced by the fixed factor metrics. The ordered
eigenvalues, counted with complex multiplicity, of
$\Delta_s=\dbar_{\End E}^{*,h_{\HE,s}}\dbar_{\End E}$ on
$L^2(X,\End_0E)$ satisfy
\begin{equation}\label{eq:main-spectrum}
\left|\frac{\lambda_j(s)}{\nu_j(s)}-1\right|\le C\sigma(s)
\quad(1\le j\le N),\qquad \lambda_{N+1}(s)\ge\gamma.
\end{equation}
If $h_{\mathrm{bal},s}$ is the pullback of $h_{0,s}$ through a
determinant-one balancing gauge and the harmonic identification of $E$,
then
\begin{equation}\label{eq:main-metric}
e^{-C\sigma(s)}h_{\mathrm{bal},s}\le h_{\HE,s}
\le e^{C\sigma(s)}h_{\mathrm{bal},s}.
\end{equation}
The constants impose no lower bound on a positive stability margin,
$\nu_1$, or any positive ratio $\nu_j/\nu_k$. No estimate of the form
$|s|=O(\sigma(s))$ is required.
\end{theorem}

The harmonic identification in this statement is fixed by removing the
exact part of the extension form, as in Lemma~\ref{lem:algebra}.
Changing the balanced representative by $K$ leaves the pulled-back metric
unchanged; the remaining common scalar is fixed by the determinant.
The constants may depend on the fixed factors, multiplicities, harmonic
spaces, and smooth parameter neighborhood. They are not asserted uniform
under degeneration of this background or growth of the multiplicities.

The theorem controls endomorphism-valued zero-forms. It is not a statement
about all degrees of the Dolbeault complex or the entire Yang--Mills
Hessian. At a nonsplit strictly semistable boundary point the original
bundle need not admit an HE metric. All occurrences of $h_{\HE,s}$ in
the theorem are on the stable side.

The example in Section~\ref{sec:example} is a fixed rank-five bundle on
$(\PP^1)^4$ with one multiplicity-two target and three scalar sources.
For genuine polarizations depending on $t,\delta>0$, it has
\begin{equation}\label{eq:intro-example}
\lambda_1=\frac54t\delta\bigl(1+O(t)+O(\delta)\bigr),\qquad
\lambda_j\asymp t\ (2\le j\le6),\qquad\lambda_7\ge\gamma.
\end{equation}
The constants are uniform for independently small $t$ and $\delta$.
In particular $\delta=e^{-1/t}$ is allowed.

\subsection{Relation to earlier work and proof structure}
The correspondence between stable bundles and HE metrics is due to
Donaldson and Uhlenbeck--Yau; we use the standard formulation and
endomorphism identities in Kobayashi~\cite{kobayashi1987}.
King~\cite{king1994} relates subrepresentation stability to moment-map
balance. Haiden--Katzarkov--Kontsevich--Pandit~\cite[\S5.1,
Theorem~5.1]{hkkp2023} give the real-weight formulation used here,
including positive masses on the multiplicity algebra.

Clarke--Tipler~\cite[Propositions~3.1--3.2]{clarketipler2024}
construct the local elliptic reduction without an abelian restriction;
their local wall theorem assumes abelian graded automorphisms.
Delloque~\cite[Proposition~4.2 and Theorems~4.8--4.9]{delloque2025}
allows nonabelian graded automorphisms and proves continuity and
quantitative convergence of HYM representatives under metric variation.
His Proposition~4.7 obtains a zero of the geometrically deformed moment
map. These results already supply nonabelian local reduction and HYM
existence. They do not identify that zero, with a relative error tending
to zero uniformly near weak directions, with a zero of the independently
specified quadratic moment map. The distinction is essential for
\eqref{eq:main-spectrum}: an additive error that is small compared with
the total extension energy may dominate the least positive eigenvalue.
All Delloque locators refer to the published version, arXiv version~4.

The repeated-factor harmonic quiver also occurs in the flow theory of
Lam~\cite[Theorems~8.2, 8.7 and Proposition~8.9]{lam2026}; the finite
iterated-logarithmic asymptotics are developed in
\cite[Theorem~5.11]{hkkp2023}. Those comparisons concern bounded metric
distance in a flow problem. The present result is stationary and gives
a multiplicative error tending to one. It extends the distinct-factor
spectral theorem of~\cite{wuebben2026uniform}; no result of that paper
is needed as an unproved step in the argument below.

The new estimate is Proposition~\ref{prop:projection}: the reduced
curvature residual, paired with any orthogonal multiplicity projection,
is bounded by total extension energy times the energy crossing that
projection. Theorem~\ref{thm:calibration} uses this estimate on spectral
projections of a Hermitian gauge logarithm to exclude zeros on the
boundary of a small ball. Brouwer degree gives the desired HE correction
without a uniform inverse for the finite Hessian. Harmonicity then makes
the Schur-complement coupling proportional to each vector's own finite
energy, yielding every relative eigenvalue estimate.
